% Generated draft: requires mathematical and source review.
\documentclass[11pt]{article}
\usepackage[margin=25mm]{geometry}
\usepackage{fontspec}
\setmainfont{Latin Modern Roman}
\usepackage{amsmath,amssymb,mathtools,unicode-math}
\setmathfont{Latin Modern Math}
\usepackage{enumitem,needspace,xurl,hyperref}
\setlength{\parindent}{0pt}
\setlength{\parskip}{5pt}
\setlength{\emergencystretch}{3em}
\newcommand{\N}{\mathbb{N}}
\newcommand{\Z}{\mathbb{Z}}
\newcommand{\Q}{\mathbb{Q}}
\newcommand{\R}{\mathbb{R}}
\newcommand{\C}{\mathbb{C}}
\newcommand{\F}{\mathbb{F}}
\newcommand{\A}{\mathbb{A}}
\newcommand{\PP}{\mathbb{P}}
\newcommand{\E}{\mathbb{E}}
\newcommand{\eps}{\varepsilon}
\newcommand{\dd}{\,\mathrm{d}}
\begin{document}
\section*{UnsolvedMath ID: 20000825}
\textbf{Source problem code:} catalogue entry
\subsection{Definitions and mathematical statement}
The original asks whether \(\operatorname{Hilb}^8(\A^4)\) is reduced and, more generally, for techniques proving reducedness of Hilbert schemes. Fix an algebraically closed field \(k\) of characteristic zero for the supplied reduction, and write \(\mathcal H=\operatorname{Hilb}^8(\A_k^4)\). Reducedness means that its local rings have no nonzero nilpotents; generic reducedness alone is weaker.

For each monomial order ideal \(B\) of size eight, containing one and closed under divisors, let \(U_B\) be the open locus where the residue classes of \(B\) form a basis of the quotient algebra. These monomial-basis opens cover \(\mathcal H\). Their coordinate rings \(A_B\) can be expressed by border-basis coefficients and the entries of commutators of the four universal multiplication matrices. The aid reports an exhaustive enumeration of 684 such bases and 49 types under coordinate permutations. Of these, 38 charts meet only one component and require proof that \(A_B\) is reduced, equivalently that a defining ideal in the polynomial presentation is radical. The other 11 have basis Hilbert function \((1,4,3)\), and meet the 25-dimensional exceptional component as well as the principal component.

For each of those 11 charts, let \(\mathfrak q_G\subset A_B\) be the prime defining the reduced exceptional component, and let \(f\in A_B\) be a suitable lift of the Salmon--Turnbull Pfaffian from the CEVV component description. The lift must have nonzero image in \(A_B/\mathfrak q_G\), namely \(f\notin\mathfrak q_G\). The aid's sufficient certificate is
\[
 A_B/(f)\text{ is reduced},\qquad
 (0:_{A_B}f)=\mathfrak q_G.
\]
The chart ring itself is then reduced. These 38 radicality checks and 11 quotient-and-colon checks constitute a finite proposed certificate; the report does not assert that they have all been completed.

The local-algebra mechanism is precise: if \(A\) is noetherian, \(\mathfrak q\) is prime, \(f\notin\mathfrak q\), \(A/(f)\) is reduced, and \((0:f)=\mathfrak q\), then \(A\) is reduced. Indeed any nilpotent lies in \((f)\), and its image in the domain \(A/\mathfrak q\) forces its coefficient to lie in \(\mathfrak q=(0:f)\). The nonvanishing hypothesis on \(f\) is essential: \(A=k[t]/(t^2)\), \(f=t\), \(\mathfrak q=(t)\) satisfies the other two conditions but is nonreduced. In the chart application the nonzero Pfaffian on the generic exceptional component is the required geometric input.

A concrete local model illustrates the need for the colon check despite apparently reduced component data:
\[
 A=k[x,y]/(x^2y,xy^2),\qquad A_{\mathrm{red}}=k[x,y]/(xy).
\]
The two reduced component quotients are \(k[x]\) and \(k[y]\), their intersection is reduced, both schemes have tangent dimension two at the origin, and the ring is generically reduced. Nevertheless \(xy\neq0\) is nilpotent. For \(f=x\), \(\mathfrak q=(y)\), one has \(A/(f)=k[y]\) but
\[
 (0:f)=(xy,y^2)\neq(y).
\]
This example is a local-algebra illustration, not a claim that this ring occurs as one of the Hilbert charts. The aid reports a worked counterexample with the same purpose but supplies no exact ring in its summary. Component dimensions, reduced component schemes, a reduced intersection, and tangent dimensions do not replace the scheme-theoretic colon equality. No answer yes or no for the whole \(\mathcal H\) is inferred before the actual chart certificates are verified.
\subsection{Short English statement}
The supplied work reduces the characteristic-zero question to a finite set of chart checks, rather than deciding reducedness. On the exceptional-component charts, reducedness of a Pfaffian quotient and an exact annihilator equality would rule out hidden nilpotent glue. The Pfaffian must also be nonzero on the exceptional component; generic and tangent data alone do not suffice.
\subsection{Sources}
\begin{itemize}
\item Local registry: permanent UnsolvedMath ID \(20000825\), exact public code \texttt{catalogue entry}; \texttt{aim-arithmetic-geometry-notes.json}, \texttt{notes[70]}.
\item AIM workshop \emph{Components of Hilbert Schemes}, July 2010, Problem 24: \url{https://aimath.org/WWN/hilbertschemes/hilbertschemes.pdf}. This supplied URL is provenance; no PDF access is claimed.
\item Supplied research summary, ``A Pfaffian-colon reducedness certificate for Hilb8(A4),'' attempt 1, cites the CEVV Pfaffian equation and a Salmon--Turnbull Pfaffian without complete author lists, paper titles, or theorem citations. It reports a 684-to-49 enumeration but supplies no chart lists or completed certificate artifacts. Original remarks and literature fields are empty.
\end{itemize}
\end{document}
